\documentclass[a4paper,11pt]{article}

\usepackage[utf8]{inputenc}
\usepackage[T1]{fontenc}
\usepackage[german]{babel}
\usepackage[a4paper,margin=2.5cm]{geometry}
\usepackage{enumitem}
\usepackage{titlesec}
\usepackage{longtable}
\usepackage{setspace}
\usepackage{parskip}

\begin{document}

% ---------------------------
% Deckblatt
% ---------------------------

\begin{titlepage}
\centering

\vspace*{3cm}

{\Huge \textbf{Meetingprotokoll}}\\[0.5cm]
{\Large Statusbericht zur Projektarbeit}\\[2cm]

\rule{12cm}{0.5pt}\\[0.5cm]

{\Large Projektarbeit}\\[0.3cm]
{\large 12. Juni 2026}

\rule{12cm}{0.5pt}

\vfill

{\large THWS / MINOVA}\\

\end{titlepage}

% ---------------------------
% Inhaltsverzeichnis
% ---------------------------

\tableofcontents
\newpage

% ---------------------------
% Inhalt
% ---------------------------

\section{Allgemeine Informationen}

\begin{tabular}{ll}
\textbf{Datum:} & 12.06.2026 \\
\textbf{Art des Meetings:} & Projektmeeting / Statusbesprechung \\
\textbf{Nächstes Meeting:} & 26.06.2026 um 09:00 Uhr \\
\end{tabular}

\vspace{1cm}

\section{Projektstatus}

Der aktuelle Stand des Projekts wurde im Meeting vorgestellt und besprochen.

Die Kartenansicht zur Anzeige der Fahrerlocation wurde erfolgreich umgesetzt und funktioniert wie vorgesehen. Außerdem wurden weitere Anpassungen für die Kundenkommunikation und die Darstellung von Datums- und Zeitangaben abgestimmt.

\subsection{Aktuell abgeschlossene bzw. bearbeitete Aufgaben}

\begin{itemize}[leftmargin=1.2cm]

\item \textbf{Systemerweiterungen}
\begin{itemize}
    \item Kartenansicht für die Fahrerlocation implementiert
    \item Anzeige von Fahrerposition und Zielort umgesetzt
    \item Erste Optimierungen der Kundenansicht durchgeführt
\end{itemize}

\item \textbf{Projektorganisation und Dokumentation}
\begin{itemize}
    \item Arbeiten an der Projektdokumentation fortgeführt
    \item Offene Fragen zur Dokumentation geklärt
\end{itemize}

\end{itemize}

\vspace{0.5cm}

\section{Besprochene Punkte}

Während des Meetings wurden folgende organisatorische und technische Themen besprochen:

\begin{itemize}[leftmargin=1.2cm]

\item Die Kartenansicht für die Fahrerlocation funktioniert erfolgreich und wurde vorgestellt.

\item Der Hinweis zur Rückmeldung soll künftig in einer benutzerfreundlicheren Form dargestellt werden, beispielsweise:

\textit{„Bitte antworten Sie spätestens vorher: 11.06.2026 9:00 Uhr“}

\item Datumsangaben sollen im Format \texttt{TT.MM.JJJJ} dargestellt werden. Formate wie \texttt{2026-06-15} sollen nicht mehr angezeigt werden.

\item Ziel ist es, bis Ende Juni alle funktionalen Anforderungen vollständig umzusetzen und ein lauffähiges Gesamtsystem präsentieren zu können.

\item Die Projektdokumentation und der Projektbericht sollen bis Ende Juli fertiggestellt werden.

\item Offene Fragen zur Projekt- und Ergebnisdokumentation wurden beantwortet.

\end{itemize}

\vspace{0.5cm}

\section{Nächste Schritte (nächste 2 Wochen)}

\begin{itemize}[leftmargin=1.2cm]
\item Fehlerbehebungen und Optimierungen
\item Weiterführung der Projektdokumentation
\end{itemize}

\vspace{0.5cm}

\section{Nächste Projektphase}

Die nächste Projektphase konzentriert sich auf:

\begin{itemize}[leftmargin=1.2cm]
\item Fertigstellung aller funktionalen Anforderungen
\item Durchführung von Systemtests
\item Erstellung der Projekt- und Ergebnisdokumentation
\end{itemize}

\vspace{0.5cm}

\section{Nächste Meilensteine}

\begin{itemize}[leftmargin=1.2cm]
\item Vollständige Umsetzung aller Anforderungen bis Ende Juni
\item Präsentation eines funktionsfähigen Gesamtsystems
\item Fertigstellung der Projektdokumentation bis Ende Juli
\end{itemize}

\vspace{1cm}

\section{Zusammenfassung}

Das Projekt befindet sich aktuell in einer fortgeschrittenen Entwicklungsphase. Die Kartenansicht zur Fahrerlocation wurde erfolgreich umgesetzt und vorgestellt. Darüber hinaus wurden weitere Verbesserungen der Kundenansicht sowie Anpassungen der Datums- und Zeitdarstellung abgestimmt.

In den kommenden Wochen liegt der Fokus insbesondere auf der Fertigstellung der verbleibenden Anforderungen, der Durchführung von Tests sowie der Erstellung der Projekt- und Ergebnisdokumentation.

\end{document}
